\capitulo{3}{Conceptos teóricos}

\section{Cloud Computing}

El cloud computing es un modelo tecnológico que permite acceder a recursos informáticos bajo demanda utilizando Internet. Este modelo reduce la necesidad de mantener infraestructuras físicas propias, permite escalar recursos dinámicamente y facilita que las organizaciones consuman capacidad de cómputo, almacenamiento, red y servicios gestionados según sus necesidades reales \cite{microsoft_azure}.

Entre las principales ventajas del cloud computing destacan:

\begin{itemize}
    \item Escalabilidad.
    \item Alta disponibilidad.
    \item Pago por uso.
    \item Reducción de costes.
    \item Facilidad de mantenimiento.
\end{itemize}

Estas ventajas son especialmente relevantes en proyectos de alcance empresarial, ya que permiten pasar de una infraestructura rígida a un modelo más flexible, automatizable y orientado a servicios. En el caso de este TFG, el uso de la nube permite desplegar un prototipo accesible para el tribunal sin distribuir máquinas virtuales ni depender de una instalación local compleja.

\subsection{Modelos de servicio}

Los servicios cloud suelen clasificarse en distintos modelos según el nivel de responsabilidad que asume el proveedor y el nivel de control que conserva el usuario. La tabla \ref{tabla:modelos_servicio_cloud} resume los modelos principales.

\begin{table}[h]
\centering
\small
\begin{tabularx}{\textwidth}{l X X}
\textbf{Modelo} & \textbf{Responsabilidad principal del proveedor} & \textbf{Ejemplo en Azure} \\
\hline
IaaS & Infraestructura básica: máquinas virtuales, red, discos y recursos de cómputo. & Azure Virtual Machines \\
PaaS & Plataforma gestionada para ejecutar aplicaciones sin administrar servidores. & Azure App Service \\
SaaS & Aplicación completa consumida como servicio por el usuario final. & Microsoft 365 \\
\end{tabularx}
\caption{Modelos principales de servicio cloud}
\label{tabla:modelos_servicio_cloud}
\end{table}

La solución desarrollada se apoya principalmente en un enfoque PaaS. Azure App Service permite ejecutar una API Flask y un portal web sin administrar directamente el sistema operativo, el balanceo interno ni la infraestructura física subyacente \cite{azure_app_service_overview}.

La tabla \ref{tabla:responsabilidad_compartida_cloud} resume cómo se reparte la responsabilidad entre el proveedor cloud y el equipo que desarrolla la aplicación. Esta distinción es importante porque el uso de servicios gestionados no elimina la responsabilidad de diseñar correctamente la aplicación, proteger los secretos y validar el despliegue.

\begin{table}[h]
\centering
\small
\begin{tabularx}{\textwidth}{l X X}
\textbf{Ámbito} & \textbf{Responsabilidad del proveedor cloud} & \textbf{Responsabilidad del proyecto} \\
\hline
Infraestructura física & Centros de datos, hardware, red física y energía. & Selección de región y uso responsable de recursos. \\
Plataforma de ejecución & Mantenimiento del servicio gestionado App Service. & Configuración de runtime, variables de entorno y publicación de la aplicación. \\
Datos & Disponibilidad del servicio de almacenamiento. & Modelo de datos, privacidad, validación y persistencia correcta. \\
Identidad y secretos & Servicios como Key Vault y Managed Identity. & Definir permisos mínimos y evitar secretos en el repositorio. \\
Aplicación & Componentes gestionados de plataforma. & Código Flask, autenticación, pruebas y control de calidad. \\
\end{tabularx}
\caption{Responsabilidad compartida en la solución cloud}
\label{tabla:responsabilidad_compartida_cloud}
\end{table}

\subsection{Modelos de despliegue}

Además de los modelos de servicio, la nube puede organizarse mediante diferentes modelos de despliegue: nube pública, privada, híbrida o multicloud. Este proyecto utiliza nube pública, ya que los recursos se despliegan en Microsoft Azure y se consumen desde una suscripción Azure for Students. Para un prototipo académico, este enfoque permite validar una arquitectura real con costes controlados y acceso remoto.

\section{Microsoft Azure}

Microsoft Azure es la plataforma cloud de Microsoft orientada al desarrollo, despliegue y administración de aplicaciones empresariales \cite{azure_architectures}. Microsoft describe Azure como una plataforma para crear, desplegar y gestionar soluciones en la nube, con servicios de cómputo, datos, inteligencia artificial, seguridad y operaciones \cite{azure_what_is}.

Los principales servicios utilizados en este proyecto son:

\begin{itemize}
    \item Azure App Service.
    \item Azure Blob Storage.
    \item Azure Key Vault.
    \item Managed Identity.
    \item Azure Monitor.
\end{itemize}

\subsection{Servicios Azure utilizados}

La tabla \ref{tabla:servicios_azure_tfg} relaciona los servicios Azure seleccionados con la responsabilidad concreta que cumplen dentro del prototipo.

\begin{table}[h]
\centering
\small
\begin{tabularx}{\textwidth}{l X X}
\textbf{Servicio} & \textbf{Uso en el TFG} & \textbf{Justificación} \\
\hline
App Service & Aloja la API Flask y el portal web. & Servicio PaaS adecuado para aplicaciones web y APIs REST. \\
Blob Storage & Guarda las solicitudes en formato JSON. & Almacenamiento sencillo, económico y apropiado para datos semiestructurados. \\
Key Vault & Almacena el token de autenticación. & Evita credenciales embebidas en código fuente. \\
Managed Identity & Permite que App Service acceda a Key Vault. & Reduce la gestión manual de secretos y claves. \\
Azure Monitor & Apoya la revisión de estado y logs. & Facilita observabilidad básica del despliegue. \\
\end{tabularx}
\caption{Servicios Azure utilizados en la solución}
\label{tabla:servicios_azure_tfg}
\end{table}

La figura \ref{fig:arquitectura-integral-azure} resume la relación entre los servicios cloud seleccionados y las evidencias del proceso de desarrollo. La suscripción de Azure contiene el grupo de recursos del TFG, donde App Service ejecuta el portal y la API Flask, Key Vault protege el secreto utilizado como token Bearer, Blob Storage conserva las solicitudes y Azure Monitor recoge telemetría. Fuera de Azure quedan las herramientas que apoyan el seguimiento y la calidad: GitHub, Zube y SonarCloud.

\begin{figure}[h]
\centering
\includegraphics[width=\textwidth]{arquitectura_integral_tfg_azure.png}
\caption{Visión integral de la arquitectura cloud y de las evidencias del TFG. Fuente: elaboración propia.}
\label{fig:arquitectura-integral-azure}
\end{figure}

Azure Blob Storage está orientado al almacenamiento de objetos y datos no estructurados, y permite acceder a los objetos mediante HTTP/HTTPS o bibliotecas cliente, incluidas bibliotecas para Python \cite{azure_blob_overview}. Azure Key Vault centraliza secretos, claves y certificados, reduciendo la probabilidad de exponer información sensible en el código \cite{azure_keyvault_overview}. Las identidades administradas permiten que un recurso de Azure obtenga una identidad propia para acceder a otros servicios sin almacenar credenciales en la aplicación \cite{azure_managed_identities}.

\section{Arquitectura basada en servicios}

Una arquitectura basada en servicios separa responsabilidades que evolucionan y se despliegan de forma distinta. En este proyecto la separación no pretende crear microservicios: busca que la interfaz, la lógica, la persistencia y la gestión de secretos tengan límites reconocibles.

Este enfoque facilita:

\begin{itemize}
    \item Escalabilidad.
    \item Mantenimiento.
    \item Reutilización.
    \item Seguridad.
\end{itemize}

La arquitectura del proyecto sigue esta idea: el portal web se encarga de la interacción con el usuario, la API Flask concentra la lógica de negocio, el clasificador aplica reglas de análisis automático, Blob Storage conserva las solicitudes y Key Vault protege el secreto usado por los endpoints protegidos. Esta separación facilita explicar el sistema, probarlo por partes y evolucionarlo en futuras versiones.

\subsection{Plano de gestión y plano de datos}

En una solución cloud conviene distinguir entre el plano de gestión y el plano de datos. El primero reúne las operaciones con las que se crean, configuran y supervisan recursos: alta del App Service, definición de variables de entorno, asignación de identidad o concesión de permisos. El segundo contiene las operaciones propias de la aplicación: registrar una solicitud, consultar su historial o guardar el documento en Blob Storage. Mezclar ambos planos haría que una petición funcional necesitara permisos de administración sobre Azure, un acoplamiento innecesario y difícil de proteger.

En el prototipo, PowerShell y Azure CLI actúan sobre el plano de gestión durante el despliegue. La API, la Logic App y el portal operan en el plano de datos durante el uso normal. Managed Identity sirve de unión controlada entre ambos: la identidad se concede desde el plano de gestión, pero App Service la utiliza durante la ejecución para obtener el secreto autorizado. Esta separación permite aplicar privilegio mínimo y explica por qué el token de usuario no debe convertirse en una credencial de administración de la suscripción.

\subsection{Persistencia documental y consistencia}

Blob Storage conserva objetos y no ofrece por sí mismo las relaciones, restricciones o consultas de un sistema gestor de bases de datos. Esa diferencia condiciona el diseño. La colección de solicitudes del prototipo se serializa como JSON y la API concentra las lecturas y escrituras; ningún cliente modifica directamente el blob. Con un único proceso y el volumen de demostración, esta solución reduce infraestructura y mantiene el formato fácil de inspeccionar. La biblioteca oficial de Azure permite ejecutar las mismas operaciones desde Python mediante HTTPS \cite{azure_blob_overview,azure_blob_python}.

La elección tiene un límite importante: dos escrituras concurrentes podrían partir de versiones distintas y sobrescribir cambios. Azure Blob Storage permite aplicar concurrencia optimista mediante ETag y cabeceras condicionales \texttt{If-Match} \cite{azure_blob_concurrency}. Una implantación con varios equipos, mayor carga o requisitos transaccionales debería utilizar ese control, separar cada solicitud en un objeto o migrar a Azure SQL o Cosmos DB. La escalabilidad del servicio cloud no corrige automáticamente un modelo de persistencia inadecuado; es necesario analizar también la granularidad y consistencia de los datos.

\section{API REST}

Una API REST permite la comunicación entre sistemas mediante peticiones HTTP.

En este proyecto se utilizan diferentes endpoints para:

\begin{itemize}
    \item Crear solicitudes operativas.
    \item Consultar solicitudes e incidencias.
    \item Obtener métricas.
\end{itemize}

\section{Seguridad y buenas prácticas cloud}

La seguridad es un aspecto transversal en soluciones cloud. En este TFG se han aplicado tres decisiones principales: uso de HTTPS en el despliegue, separación de secretos mediante Key Vault y autenticación Bearer para las operaciones de consulta. Estas medidas no convierten el prototipo en una solución corporativa completa, pero sí demuestran buenas prácticas básicas de diseño seguro.

\begin{sloppypar}
Además, el diseño se ha alineado con principios del Azure Well-Architected Framework, que propone pilares como seguridad, fiabilidad, eficiencia de rendimiento, optimización de costes y excelencia operativa \cite{azure_well_architected_framework}. En el contexto del proyecto, estos principios se reflejan en el uso de servicios gestionados, despliegue reproducible, recursos de bajo coste y documentación de evidencias.
\end{sloppypar}

\subsection{Observabilidad frente a monitorización}

La monitorización responde a preguntas conocidas, por ejemplo cuántas peticiones recibe el servicio o si se producen respuestas con error. La observabilidad busca además poder investigar comportamientos no previstos a partir de métricas, registros y trazas. Application Insights y Azure Monitor cubren esa necesidad básica en el proyecto: permiten comprobar actividad del App Service, localizar errores y relacionar el estado técnico con las pruebas funcionales \cite{azure_monitor}.

El endpoint \texttt{/health} complementa estas herramientas, pero no las sustituye. Su respuesta confirma que la aplicación está disponible y muestra el modo de almacenamiento; no demuestra que todos los procesos empresariales funcionen. Por esa razón la validación combina tres niveles: salud técnica, pruebas automáticas del código y recorrido completo sobre el despliegue. Esta combinación evita considerar correcto un sistema únicamente porque su proceso web responde.

\section{DevOps y automatización}

DevOps es una metodología orientada a integrar desarrollo y operaciones con el objetivo de automatizar procesos y mejorar la calidad del software \cite{bojja2020_devops}.

De las prácticas asociadas a DevOps, el proyecto utiliza las siguientes:

\begin{itemize}
    \item Control de versiones y commits vinculados al seguimiento.
    \item Despliegue repetible mediante PowerShell y Azure CLI.
    \item Verificación automática de los endpoints publicados.
    \item Revisión externa de calidad con SonarCloud.
\end{itemize}

La versión final no mantiene un pipeline de integración continua. Los scripts de despliegue y verificación cubren la necesidad real del proyecto: repetir la publicación desde el equipo de desarrollo y dejar visibles las órdenes ejecutadas sobre Azure. Esta diferencia respecto al plan inicial se documenta en el seguimiento de sprints.
